\chapter{Estado del Arte}

\section{Soluciones Comerciales Vigentes}

El mercado de herramientas de gestión y optimización cloud ha experimentado un crecimiento significativo en los últimos años. Entre las soluciones comerciales más destacadas se encuentran:

\subsection{Amazon CloudWatch y AWS Cost Explorer}

Amazon CloudWatch es un servicio de monitoreo y observabilidad proporcionado por Amazon Web Services que permite recopilar y supervisar métricas, registros (logs) y eventos generados por aplicaciones, servicios e infraestructura desplegada en la nube. La plataforma ofrece funcionalidades de visualización mediante paneles de control, generación de alertas y análisis del comportamiento de los recursos con el objetivo de facilitar la supervisión operativa y la detección de incidentes \parencite{AWS2026}.

Complementariamente, AWS Cost Explorer es una herramienta orientada al análisis y administración de costos que permite visualizar el consumo de servicios cloud, identificar tendencias de gasto, generar reportes y realizar estimaciones de costos futuros. Este servicio proporciona recomendaciones relacionadas con la optimización del consumo dentro del ecosistema AWS \parencite{AWS2026}.

AWS Trusted Advisor es una herramienta que analiza la configuración y utilización de los recursos desplegados en AWS con el objetivo de proporcionar recomendaciones relacionadas con optimización de costos, rendimiento, seguridad, tolerancia a fallos y límites de servicio. Para ello, evalúa la infraestructura mediante un conjunto de verificaciones predefinidas (checks) y genera sugerencias orientadas a identificar recursos subutilizados, configuraciones mejorables o posibles riesgos operativos \parencite{AWS2026}.

En conjunto, Amazon CloudWatch, AWS Cost Explorer y AWS Trusted Advisor proporcionan capacidades de monitoreo, análisis de costos y generación de recomendaciones para la administración de recursos dentro del ecosistema AWS. Mientras CloudWatch permite supervisar métricas, registros y eventos operativos, Cost Explorer facilita el análisis del gasto asociado al consumo de servicios cloud y Trusted Advisor complementa estas funcionalidades mediante verificaciones y recomendaciones predefinidas sobre la infraestructura. Sin embargo, las recomendaciones disponibles dependen de los controles implementados por la plataforma y se encuentran asociadas específicamente a los servicios y recursos soportados por AWS \parencite{AWS2026}.

\subsection{Google Cloud Operations Suite}

Google Cloud Operations Suite es el conjunto de herramientas de monitoreo, registro y diagnóstico de Google Cloud destinado a supervisar aplicaciones, servicios e infraestructura desplegados en la nube. La plataforma integra funcionalidades de monitoreo de métricas, gestión de registros (logging), trazabilidad distribuida (tracing), generación de alertas y análisis del rendimiento de sistemas, permitiendo obtener visibilidad sobre el comportamiento operativo de los recursos y facilitar la identificación de incidentes \parencite{GoogleCloud2026}.

La suite proporciona capacidades de observabilidad orientadas a recopilar y correlacionar información proveniente de distintos componentes de la infraestructura, contribuyendo al análisis del estado de los sistemas y al diagnóstico de problemas de rendimiento y disponibilidad. Entre sus funcionalidades se incluyen paneles de visualización, monitoreo de aplicaciones, análisis de eventos y herramientas para la administración de alertas \parencite{GoogleCloud2026}.

Google Cloud Active Assist es un conjunto de herramientas de análisis y recomendación que utiliza datos operativos y técnicas de aprendizaje automático para ayudar a los usuarios a optimizar recursos, mejorar el rendimiento, reforzar la seguridad y administrar costos dentro de Google Cloud. La plataforma analiza continuamente el comportamiento de la infraestructura y genera recomendaciones sobre distintos servicios, incluyendo máquinas virtuales, bases de datos, cuotas de recursos e identidades y accesos, con el objetivo de favorecer una utilización más eficiente de los recursos disponibles. Entre sus funcionalidades se encuentran la identificación de recursos ociosos o sobredimensionados, la detección de configuraciones que podrían afectar el rendimiento de los sistemas y la generación de sugerencias orientadas a la reducción de costos operativos. Estas recomendaciones son presentadas al usuario mediante interfaces integradas dentro de Google Cloud, facilitando la toma de decisiones sobre la infraestructura desplegada \parencite{GoogleCloud2026}.

En conjunto, Google Cloud Operations Suite y Active Assist proporcionan capacidades avanzadas de observabilidad, monitoreo y análisis para recursos desplegados en Google Cloud. Mientras Operations Suite se orienta a la recopilación y visualización de métricas, registros y trazas, Active Assist complementa estas funcionalidades mediante un conjunto de recomendadores (recommenders) especializados que analizan el comportamiento de determinados servicios y generan sugerencias de optimización relacionadas con costos, rendimiento, utilización de recursos y configuración de infraestructura. Sin embargo, las recomendaciones disponibles dependen de los tipos de recursos y recomendadores implementados por la plataforma, y requieren la existencia de un historial de uso suficiente para que el sistema pueda analizar patrones de comportamiento y generar insights relevantes \parencite{GoogleCloud2026}.

\subsection{Azure Monitor y Azure Advisor}

Azure Monitor es la plataforma de monitoreo y observabilidad de Microsoft Azure, diseñada para recopilar, almacenar, analizar y visualizar información proveniente de aplicaciones, servicios e infraestructura desplegados en la nube. La herramienta permite supervisar el comportamiento de recursos tanto a nivel de infraestructura como de aplicaciones, proporcionando información sobre rendimiento, disponibilidad y utilización de recursos mediante la recopilación continua de métricas, registros y eventos operativos \parencite{Microsoft2026}.

La plataforma integra diferentes mecanismos de observabilidad que permiten centralizar información generada por máquinas virtuales, bases de datos, redes, contenedores y aplicaciones. Entre sus funcionalidades se encuentran el monitoreo en tiempo real, la generación de alertas automáticas, la visualización mediante paneles de control personalizables y el análisis histórico de métricas. Además, Azure Monitor incorpora herramientas para la gestión de registros (Log Analytics), trazabilidad distribuida de aplicaciones y análisis de eventos, facilitando el diagnóstico de incidentes y la identificación de problemas de rendimiento dentro de arquitecturas distribuidas \parencite{Microsoft2026}.

Azure Advisor es un servicio de recomendaciones que analiza la configuración y utilización de los recursos desplegados en Microsoft Azure con el objetivo de identificar oportunidades de mejora dentro de la infraestructura. A partir de la evaluación continua de los servicios utilizados, la herramienta genera recomendaciones personalizadas orientadas a optimizar costos, incrementar el rendimiento de las aplicaciones, mejorar la disponibilidad de los sistemas, reforzar la seguridad y promover buenas prácticas operativas \parencite{Microsoft2026}.

Las recomendaciones son presentadas al usuario a través del portal de Azure y se encuentran organizadas en distintas categorías, entre las que se incluyen confiabilidad, seguridad, rendimiento, excelencia operativa y optimización de costos. Estas sugerencias pueden abarcar acciones como la identificación de recursos subutilizados, el ajuste de configuraciones, la adopción de mecanismos de alta disponibilidad o la implementación de medidas de seguridad recomendadas por Microsoft. De esta manera, Azure Advisor actúa como un mecanismo de asistencia para la toma de decisiones relacionadas con la administración de infraestructura cloud \parencite{Microsoft2026}.

En conjunto, Azure Monitor y Azure Advisor proporcionan capacidades integradas de observabilidad, monitoreo y generación de recomendaciones para la administración de recursos dentro del ecosistema Microsoft Azure. Mientras Azure Monitor se orienta a la recopilación, almacenamiento y análisis de información operativa proveniente de la infraestructura y las aplicaciones, Azure Advisor utiliza dicha información junto con análisis propios de la plataforma para generar recomendaciones orientadas a mejorar el rendimiento, la seguridad, la disponibilidad y la eficiencia económica de los recursos desplegados. Sin embargo, las recomendaciones disponibles dependen de los mecanismos de análisis implementados por Microsoft y se encuentran asociadas específicamente a los servicios soportados dentro del entorno Azure \parencite{Microsoft2026}.

\subsection{Huawei Cloud Eye y Huawei Cloud Cost Center}

Huawei Cloud Eye es el servicio de monitoreo de Huawei Cloud destinado a supervisar el estado y comportamiento de los recursos desplegados dentro de la plataforma. La herramienta permite recopilar métricas de rendimiento provenientes de distintos servicios cloud, incluyendo instancias de cómputo, almacenamiento, redes y bases de datos, proporcionando visibilidad sobre la utilización y disponibilidad de la infraestructura \parencite{HuaweiCloud2026}.

Entre sus principales funcionalidades se encuentran la recolección continua de métricas, la visualización de información mediante gráficos y paneales de control, la configuración de alarmas y notificaciones automáticas, y el análisis histórico del comportamiento de los recursos. Estas capacidades permiten a los administradores identificar problemas operativos, supervisar el consumo de recursos y realizar tareas de seguimiento sobre el rendimiento de los servicios desplegados en la nube \parencite{HuaweiCloud2026}.

A su vez, Cloud Eye se integra con otros servicios del ecosistema Huawei Cloud para centralizar la supervisión de la infraestructura y facilitar la administración operativa del entorno. Es decir, la plataforma proporciona información relevante para la detección de incidentes y el monitoreo continuo de sistemas cloud, constituyendo uno de los principales mecanismos de observabilidad disponibles dentro del proveedor. Huawei Cloud Cost Center es una herramienta orientada a la gestión financiera de recursos cloud que permite visualizar, analizar y controlar los costos asociados al consumo de servicios dentro de Huawei Cloud. La plataforma proporciona información sobre gastos históricos, tendencias de consumo y distribución de costos entre distintos recursos y servicios utilizados por la organización \parencite{HuaweiCloud2026}.

Entre sus funcionalidades se incluyen el seguimiento del gasto cloud, la generación de reportes de costos, la definición de presupuestos y el monitoreo del consumo de recursos desde una perspectiva financiera. Estas capacidades permiten obtener mayor visibilidad sobre el impacto económico de la infraestructura desplegada y facilitan la adopción de prácticas de control y optimización del gasto tecnológico \parencite{HuaweiCloud2026}.

En conjunto, Huawei Cloud Eye y Huawei Cloud Cost Center proporcionan capacidades de monitoreo operativo y gestión financiera para la administración de recursos dentro del ecosistema Huawei Cloud. Mientras Cloud Eye se enfoca en la recopilación y visualización de métricas relacionadas con el rendimiento y disponibilidad de la infraestructura, Cost Center permite analizar los costos asociados a la utilización de los servicios cloud. Sin embargo, las funcionalidades ofrecidas por estas herramientas se encuentran principalmente orientadas al monitoreo y análisis de información, sin incorporar un mecanismo integral de detección de anomalías basado en aprendizaje automático no supervisado que permita identificar patrones inusuales de utilización de recursos y generar recomendaciones específicas de optimización sobre la infraestructura desplegada \parencite{HuaweiCloud2026}.

\subsection{VMware CloudHealth}

VMware CloudHealth es una plataforma de gestión y optimización de entornos cloud diseñada para proporcionar visibilidad sobre el consumo de recursos, los costos operativos y el cumplimiento de políticas dentro de infraestructuras desplegadas en múltiples proveedores cloud. La herramienta recopila información proveniente de servicios de computación, almacenamiento, redes y otros recursos cloud con el objetivo de centralizar su análisis y facilitar la administración de entornos complejos \parencite{VMware2023}.

Entre sus principales funcionalidades se encuentran el monitoreo del consumo de recursos, el análisis de costos cloud, la asignación de gastos entre equipos o proyectos y la generación de reportes orientados a la adopción de prácticas FinOps. La plataforma permite establecer políticas de gobernanza y supervisar el cumplimiento de configuraciones definidas por la organización, proporcionando una visión integral tanto desde una perspectiva técnica como financiera \parencite{VMware2023}.

CloudHealth incorpora capacidades de optimización mediante la generación de recomendaciones destinadas a mejorar la eficiencia en la utilización de recursos y reducir costos operativos. Estas recomendaciones se basan en el análisis de métricas históricas de consumo y en la identificación de patrones asociados a recursos infrautilizados, instancias sobredimensionadas, almacenamiento sin uso y otras oportunidades de optimización. Por esto, la plataforma asiste a los administradores en la toma de decisiones relacionadas con la gestión y el dimensionamiento de la infraestructura cloud \parencite{VMware2023}.

En conjunto, las funcionalidades de monitoreo, análisis financiero y optimización de VMware CloudHealth proporcionan una solución orientada a la gestión eficiente de recursos cloud dentro de entornos empresariales. Sin embargo, sus capacidades de recomendación se encuentran principalmente enfocadas en la optimización financiera y la gobernanza de infraestructura, mientras que el análisis realizado por la plataforma se basa en mecanismos y criterios definidos por el proveedor. Su adopción se encuentra orientada principalmente a los ecosistemas cloud de mayor presencia en el mercado, como AWS, Microsoft Azure y Google Cloud Platform \parencite{VMware2023}.

\subsection{Datadog}

Datadog es una plataforma de monitoreo y seguridad basada en la nube que proporciona observabilidad integral para aplicaciones, infraestructura y servicios en entornos cloud modernos. La plataforma integra capacidades de monitoreo de infraestructura, aplicaciones, logs, seguridad y experiencia digital en una solución unificada, permitiendo a las organizaciones obtener visibilidad completa sobre sus stacks tecnológicos a cualquier escala \parencite{Datadog2026}.

Entre sus principales funcionalidades se encuentran el monitoreo de infraestructura con visualización detallada de hosts, contenedores y recursos cloud; Application Performance Monitoring (APM) con trazabilidad distribuida y perfiles de código; gestión de logs con análisis y correlación automática; monitoreo de experiencia digital mediante Real User Monitoring (RUM) y Synthetic Monitoring; y capacidades de seguridad cloud que incluyen Cloud Security Posture Management (CSPM) y detección de amenazas \parencite{Datadog2026}.

Datadog incorpora funcionalidades avanzadas de gestión de costos cloud mediante Cloud Cost Management, que permite visualizar y optimizar el gasto en múltiples proveedores cloud incluyendo AWS, Azure y Google Cloud. La plataforma proporciona recomendaciones de optimización de costos basadas en el análisis de patrones de uso, identificación de recursos infrautilizados y sugerencias de dimensionamiento adecuado de instancias \parencite{Datadog2026}.

La plataforma destaca por su integración nativa con más de 800 tecnologías y servicios, incluyendo Kubernetes, Docker, AWS, Azure, Google Cloud, bases de datos, sistemas de mensajería y herramientas de CI/CD. Además, incorpora capacidades de inteligencia artificial mediante Bits AI, que proporciona asistentes impulsados por IA para investigación, análisis y remediación automática de incidentes \parencite{Datadog2026}.

En conjunto, Datadog proporciona una solución integral de observabilidad que combina monitoreo de infraestructura, aplicaciones, logs, seguridad y experiencia digital en una plataforma unificada. Su enfoque multi-cloud y las capacidades de correlación automática entre diferentes fuentes de telemetría permiten obtener una visión completa del comportamiento de sistemas distribuidos. Sin embargo, su modelo de pricing basado en volumen de datos ingeridos puede resultar costoso para organizaciones con grandes volúmenes de telemetría, y las recomendaciones de optimización se basan principalmente en reglas predefinidas y análisis estadísticos \parencite{Datadog2026}.

\subsection{Dynatrace}

Dynatrace es una plataforma de observabilidad impulsada por inteligencia artificial diseñada para proporcionar visibilidad completa sobre aplicaciones, infraestructura cloud y experiencia digital. La plataforma utiliza un enfoque automático de instrumentación mediante su agente OneAgent, que despliega automáticamente la instrumentación para aplicaciones, infraestructura y servicios sin requerir configuración manual extensiva \parencite{Dynatrace2026}.

Entre sus principales funcionalidades se encuentran el monitoreo de aplicaciones con trazabilidad distribuida automática y análisis de rendimiento de código; observabilidad de infraestructura para entornos multi-cloud e híbridos; análisis de logs con capacidades de búsqueda y correlación; monitoreo de experiencia digital mediante Real User Monitoring y Synthetic Monitoring; y observabilidad de inteligencia artificial para aplicaciones que utilizan LLMs y agentes de IA \parencite{Dynatrace2026}.

Dynatrace incorpora Davis AI, un motor de inteligencia artificial determinista que proporciona detección automática de anomalías, análisis de causa raíz y predicción de problemas. A diferencia de enfoques basados en machine learning probabilístico, Davis AI utiliza modelos causales deterministas que permiten identificar con precisión la causa raíz de problemas y generar recomendaciones accionables sin generar falsos positivos \parencite{Dynatrace2026}.

La plataforma proporciona capacidades avanzadas de automatización mediante Dynatrace Intelligence, que permite predecir problemas, detectar anomalías y coordinar acciones automatizadas a través de plataformas cloud, herramientas de desarrollo y sistemas de gestión de servicios de TI. Además, incorpora Grail, un data lakehouse que permite almacenar y analizar datos de telemetría en contexto, facilitando consultas complejas y análisis avanzados \parencite{Dynatrace2026}.

En conjunto, Dynatrace proporciona una solución de observabilidad con enfoque en automatización e inteligencia artificial determinista. Su capacidad de instrumentación automática reduce la complejidad de configuración, mientras que Davis AI proporciona análisis causal preciso para identificación de problemas. Sin embargo, su enfoque propietario y el modelo de pricing pueden representar una inversión significativa para organizaciones, y las capacidades de optimización de costos cloud son menos desarrolladas comparadas con soluciones especializadas en FinOps \parencite{Dynatrace2026}.

\subsection{New Relic}

New Relic es una plataforma de observabilidad inteligente que proporciona visibilidad completa sobre el stack tecnológico, desde infraestructura hasta experiencia del usuario. La plataforma integra monitoreo de aplicaciones, infraestructura, logs, experiencia digital y seguridad en una solución unificada, permitiendo correlacionar datos de diferentes fuentes para obtener una visión integral del comportamiento de los sistemas \parencite{NewRelic2026}.

Entre sus principales funcionalidades se encuentran Application Performance Monitoring (APM) con trazabilidad distribuida y análisis de transacciones; monitoreo de infraestructura para entornos cloud, Kubernetes y sistemas híbridos; gestión de logs con análisis inteligente y correlación automática; monitoreo de experiencia digital mediante Browser Monitoring, Mobile Monitoring y Synthetic Monitoring; y capacidades de seguridad con gestión de vulnerabilidades \parencite{NewRelic2026}.

New Relic incorpora capacidades avanzadas de inteligencia artificial mediante New Relic AI y AIOps, que proporcionan detección de anomalías, predicción de incidentes y automatización de respuestas. La plataforma incluye SRE Agent, un agente impulsado por IA que automatiza la investigación y remediación de incidentes, reduciendo el tiempo de resolución y la carga operativa sobre los equipos de ingeniería \parencite{NewRelic2026}.

La plataforma proporciona Cloud Cost Intelligence, una funcionalidad que permite visualizar y optimizar costos en entornos multi-cloud y Kubernetes. Esta herramienta proporciona visibilidad en tiempo real sobre el gasto cloud, identifica oportunidades de optimización y permite correlacionar costos con el rendimiento y utilización de recursos. Además, New Relic soporta OpenTelemetry como estándar abierto para la recopilación de telemetría \parencite{NewRelic2026}.

En conjunto, New Relic proporciona una plataforma de observabilidad con enfoque en inteligencia artificial y automatización. Su modelo de pricing basado en datos ingeridos ofrece flexibilidad para organizaciones de diferentes tamaños, y las capacidades de correlación automática facilitan el análisis de sistemas distribuidos. Sin embargo, la profundidad de las recomendaciones de optimización de costos puede ser limitada comparada con soluciones especializadas, y la configuración inicial puede requerir esfuerzo significativo para obtener observabilidad completa \parencite{NewRelic2026}.

\section{Soluciones de Código Abierto}

Además de las soluciones comerciales, existen herramientas de código abierto que ofrecen capacidades de monitoreo y observabilidad para entornos cloud:

\subsection{Prometheus}

Prometheus es un sistema de monitoreo y alertamiento de código abierto que recopila métricas de aplicaciones y servicios mediante un modelo de extracción (pull model). La herramienta utiliza un modelo de datos multidimensional con identificadores de métricas y etiquetas, permitiendo consultas flexibles mediante su lenguaje PromQL. Prometheus es ampliamente adoptado en entornos Kubernetes y proporciona capacidades de almacenamiento de series temporales, generación de alertas y visualización mediante integración con herramientas como Grafana \parencite{Prometheus2026}.

\subsection{Grafana}

Grafana es una plataforma de visualización y análisis de código abierto que permite crear paneles de control interactivos para métricas provenientes de múltiples fuentes de datos. La herramienta soporta integración con Prometheus, InfluxDB, Elasticsearch y otros sistemas de almacenamiento de métricas, proporcionando capacidades de visualización, alertamiento y análisis de datos operativos \parencite{Grafana2026}.

\subsection{OpenTelemetry}

OpenTelemetry es un proyecto de código abierto que proporciona un conjunto de APIs, SDKs y herramientas para la instrumentación, generación, recopilación y exportación de datos de telemetría (métricas, registros y trazas). El proyecto busca proporcionar un estándar vendor-neutral para la observabilidad, permitiendo a las organizaciones recopilar datos de telemetría independientemente del proveedor cloud o del backend de análisis utilizado \parencite{OpenTelemetry2026}.

\section{Análisis Comparativo}

El análisis realizado muestra que las principales soluciones comerciales de gestión cloud, como AWS CloudWatch y Trusted Advisor, Google Cloud Operations Suite y Active Assist, Azure Monitor y Advisor, Huawei Cloud Eye con Cost Center, Datadog, Dynatrace y New Relic ofrecen capacidades avanzadas de monitoreo, observabilidad y optimización de recursos. Estas plataformas permiten supervisar métricas operativas, visualizar el consumo de infraestructura y generar recomendaciones orientadas a reducir costos o mejorar el rendimiento de los sistemas.

No obstante, también se identifican limitaciones comunes que restringen su alcance frente a las necesidades planteadas en el presente proyecto. En primer lugar, existe una fuerte dependencia del proveedor cloud en las herramientas nativas, ya que se encuentran diseñadas para operar principalmente dentro de sus propios ecosistemas, dificultando la portabilidad y la adaptación a otros entornos. Esta situación resulta especialmente relevante en el caso de Huawei Cloud, donde la oferta de herramientas avanzadas de optimización es más limitada que en AWS, Azure o Google Cloud.

En segundo lugar, las recomendaciones generadas por estas plataformas se basan mayormente en reglas, verificaciones o recommenders predefinidos por el proveedor. Si bien algunos servicios, como Active Assist y Dynatrace Davis AI, incorporan mecanismos de análisis más sofisticados, el usuario no dispone de control ni visibilidad sobre el modelo utilizado ni puede adaptar fácilmente el comportamiento de detección a patrones específicos de su organización. En contraste, el presente proyecto propone un enfoque explícito y documentado basado en Isolation Forest, un algoritmo de aprendizaje automático no supervisado orientado a identificar anomalías directamente a partir de las métricas operativas recolectadas.

Asimismo, las soluciones comerciales analizadas presentan una escasa integración con prácticas de Infraestructura como Código (IaC). Aunque permiten detectar oportunidades de optimización, la implementación de los cambios recomendados suele requerir intervención manual o herramientas externas. La propuesta de esta tesis incorpora un módulo integrado con Terraform, permitiendo automatizar el despliegue y la administración de infraestructura desde la misma plataforma.

Otra diferencia importante radica en el enfoque del análisis. Las herramientas existentes se orientan principalmente a un modelo reactivo, donde los problemas se detectan una vez que ya se manifiestan en la infraestructura o en el consumo. La solución propuesta busca avanzar hacia un enfoque más proactivo, utilizando aprendizaje automático no supervisado para detectar comportamientos inusuales y generar recomendaciones antes de que las ineficiencias se consoliden como problemas operativos o económicos.

En conjunto, estas observaciones evidencian una brecha entre las capacidades actuales del mercado y la necesidad de contar con una plataforma que integre monitoreo continuo, detección inteligente de anomalías, recomendaciones adaptadas al contexto operacional y automatización de infraestructura dentro de un entorno Huawei Cloud. El proyecto de tesis se posiciona precisamente en esa intersección, proponiendo una solución unificada que combine monitoreo, observabilidad, análisis basado en Machine Learning y despliegue automatizado mediante IaC.
